% ============ 第3章 §3.3 A(T)中的谱综合(原书 pp.237-249) ============
\section{\texorpdfstring{$A(\mathbf{T})$ 中的谱综合}{A(T) 中的谱综合}}\label{sec:3.3}
\index{Malliavin 定理}
\index{Kronecker 集}
\index{Kronecker 定理}
\index{Cantor 集}
\index{在 Q 上线性无关}


这里,首先叙述谱综合集的若干例子,后半部分叙述证明 $A(\mathbf{T})$ 谱综合不可能的 Malliavin 的结果 \cite{ch3:34}.其后,叙述作为稀疏集(thin set)的 Kronecker 集.

\subsection{谱综合集的例子}\label{sec:3.3.1}

\begin{theorem}\label{thm:3.3.1}
$\mathbf{T}$ 上的一点是 $A(\mathbf{T})$ 的谱综合集.
\end{theorem}

\begin{proof}
考虑平移,只需证明 $\{0\} \subset \mathbf{T} = [-\pi, \pi)$ 是谱综合集.因为,设 $\{0\}$ 是谱综合集.取 $x_{0} \in \mathbf{T}$,$f \in A(\mathbf{T})$,$f(x_{0}) = 0$.令 $f_{-x_{0}}(x) = f(x + x_{0})$,则 $f_{-x_{0}} \in A(\mathbf{T})$,$f_{-x_{0}}(0) = f(x_{0}) = 0$.因 $\{0\}$ 是谱综合集,对任意 $\varepsilon > 0$ 存在 $g \in A(\mathbf{T})$ 使得 $g$ 在原点某邻域内为 $0$ 且 $\|f_{-x_{0}} - g\|_{A(\mathbf{T})} < \varepsilon$.令 $h(x) = g(x - x_{0})$,则 $h \in A(\mathbf{T})$,$h$ 在 $x_{0}$ 的某邻域内为 $0$.这里,由 $f(x + x_{0}) - h(x + x_{0}) = f_{-x_{0}}(x) - g(x)$,$\|f - h\|_{A(\mathbf{T})} = \|f_{-x_{0}} - g\|_{A(\mathbf{T})} < \varepsilon$.于是下面证明 $\{0\}$ 是 $A(\mathbf{T})$ 的谱综合集.设 $f(0) = 0$,$f \in A(\mathbf{T})$.现在,令 $0 < \varepsilon < 1/2$,
\begin{equation*}
\Delta_{\varepsilon}(x) = \max\left\{0, 1 - \frac{|x|}{\varepsilon}\right\}, \quad V_{\varepsilon} = 2\Delta_{2\varepsilon} - \Delta_{\varepsilon},
\end{equation*}
则
\begin{equation*}
\lim_{\varepsilon \to 0}\|V_{\varepsilon}f\|_{A(\mathbf{T})} = 0.
\end{equation*}
事实上,对连续可微函数 $f \in C^{1}(\mathbf{T})$,
\begin{equation*}
\|V_{\varepsilon}f\|_{A(\mathbf{T})} = \sum_{n \in \mathbf{Z}}\left|\widehat{V_{\varepsilon}f}(n)\right| = \left|\widehat{V_{\varepsilon}f}(0)\right| + \sum_{n \neq 0}\left|\widehat{V_{\varepsilon}f}(n)\right|.
\end{equation*}
这里考虑
\begin{equation*}
\left|\widehat{V_{\varepsilon}f}(0)\right| = \left|\frac{1}{2\pi}\int_{-\pi}^{\pi}V_{\varepsilon}(x)f(x)\dd x\right| = \left|\frac{1}{2\pi}\int_{-2\varepsilon}^{2\varepsilon}V_{\varepsilon}(x)f(x)\dd x\right|.
\end{equation*}
由
\begin{equation*}
\int_{-2\varepsilon}^{2\varepsilon}V_{\varepsilon}(x)f(x)\frac{\dd x}{2\pi} = \left[x(V_{\varepsilon}f)(x)\frac{1}{2\pi}\right]_{-2\varepsilon}^{2\varepsilon} - \frac{1}{2\pi}\int_{-2\varepsilon}^{2\varepsilon}x(V_{\varepsilon}f)'(x)\dd x
\end{equation*}
对右边用 Schwarz 不等式,则
\begin{align*}
\left|\int_{-2\varepsilon}^{2\varepsilon}V_{\varepsilon}(x)f(x)\frac{\dd x}{2\pi}\right| & \leq \left|\int_{-2\varepsilon}^{2\varepsilon}x(V_{\varepsilon}f)'(x)\frac{\dd x}{2\pi}\right| \\
& \leq \left(\int_{-2\varepsilon}^{2\varepsilon}x^{2}\frac{\dd x}{2\pi}\right)^{1/2}\left(\int_{-\pi}^{\pi}\left|(V_{\varepsilon}f)'(x)\right|^{2}\frac{\dd x}{2\pi}\right)^{1/2} \\
& = C\varepsilon^{3/2}\left(\int_{-\pi}^{\pi}\left|(V_{\varepsilon}f)'(x)\right|^{2}\frac{\dd x}{2\pi}\right)^{1/2}
\end{align*}
其次,当 $n \neq 0$ 时,
\begin{align*}
\widehat{V_{\varepsilon}f}(n) & = \frac{1}{2\pi}\int_{-\pi}^{\pi}(V_{\varepsilon}f)(x)e^{-inx}\dd x \\
& = \left[-\frac{e^{-inx}}{2\pi in}(V_{\varepsilon}f)(x)\right]_{-\pi}^{\pi} + \frac{1}{2\pi in}\int_{-\pi}^{\pi}(V_{\varepsilon}f)'(x)e^{-inx}\dd x
\end{align*}
故
\begin{equation*}
\widehat{V_{\varepsilon}f}(n) = \frac{1}{2\pi in}\int_{-\pi}^{\pi}(V_{\varepsilon}f)'(x)e^{inx}\dd x.
\end{equation*}
又,
\begin{align*}
\sum_{n \neq 0}\left|\widehat{V_{\varepsilon}f}(n)\right| & = \sum_{n \neq 0}\frac{1}{n}\left|\int_{-\pi}^{\pi}(V_{\varepsilon}f)'(x)e^{-inx}\frac{\dd x}{2\pi}\right| \\
& \leq \left(\sum_{n \neq 0}\frac{1}{n^{2}}\right)^{1/2}\left(\sum_{n \neq 0}\left|\widehat{(V_{\varepsilon}f)'}(n)\right|^{2}\right)^{1/2} \\
& \leq C\left(\frac{1}{2\pi}\int_{-\pi}^{\pi}\left|(V_{\varepsilon}f)'(x)\right|^{2}\dd x\right)^{1/2}.
\end{align*}
这里,
\begin{equation*}
\left(\int_{-\pi}^{\pi}\left|(V_{\varepsilon}f)'(x)\right|^{2}\frac{\dd x}{2\pi}\right)^{1/2} \leq C\varepsilon^{1/2}.
\end{equation*}
事实上,在几乎所有点上 $(V_{\varepsilon}f)'(x) = V_{\varepsilon}'(x)f(x) + V_{\varepsilon}(x)f'(x)$,由
\begin{equation*}
V_{\varepsilon}'(x) = \left\{\begin{array}{l l}
\frac{1}{2\varepsilon} & (-2\varepsilon < x < -\varepsilon) \\
-\frac{1}{2\varepsilon} & (\varepsilon < x < 2\varepsilon) \\
0 & (|x| \leq \varepsilon,\ |x| \geq 2\varepsilon)
\end{array}\right.
\end{equation*}
与 $f(0) = 0$ 即得.综上,
\begin{equation*}
\lim_{\varepsilon \to 0}\|V_{\varepsilon}f\|_{A(\mathbf{T})} = 0.
\end{equation*}
对一般的 $f \in A(\mathbf{T})$,由连续可微函数全体在 $A(\mathbf{T})$ 中稠密与 $\|V_{\varepsilon}\|_{A(\mathbf{T})} \leq 3$ 即得.$\|V_{\varepsilon}\|_{A(\mathbf{T})} \leq 3$ 由
\begin{equation*}
\widehat{\Delta_{\varepsilon}}(n) = \frac{2}{\pi} \cdot \frac{\sin^{2}\left(\frac{\varepsilon n}{2}\right)}{\varepsilon n^{2}} \geq 0, \quad \|\Delta_{\varepsilon}\|_{A(\mathbf{T})} = \Delta_{\varepsilon}(0) = 1
\end{equation*}
可知.于是设 $f \in A(\mathbf{T})$,$f(0) = 0$.令 $g_{\varepsilon} = (1 - V_{\varepsilon})f$,则 $g_{\varepsilon}$ 在原点某邻域内为 $0$,且 $g_{\varepsilon} \in A(\mathbf{T})$.又,
\begin{equation*}
\|g_{\varepsilon} - f\|_{A(\mathbf{T})} = \|V_{\varepsilon}f\|_{A(\mathbf{T})} \to 0 \quad (\varepsilon \to 0)
\end{equation*}
故
\begin{equation*}
\lim_{\varepsilon \to 0}\|g_{\varepsilon} - f\|_{A(\mathbf{T})} = 0.
\end{equation*}
因此证明了 $\{0\}$ 是 $A(\mathbf{T})$ 的谱综合集.
\end{proof}

\begin{corollary}\label{cor:3.3.2}
闭区间 $[a, b] \subset \mathbf{T}$ 是 $A(\mathbf{T})$ 的谱综合集.
\end{corollary}

\begin{proof}
不妨设 $[a, b] \subset (-\pi, \pi)$.令 $V_{\varepsilon}^{a}(x) = V_{\varepsilon}(x - a)$,$V_{\varepsilon}^{b}(x) = V_{\varepsilon}(x - b)$.设 $f \in A(\mathbf{T})$ 满足 $f(x) = 0$($x \in [a, b]$).令 $g_{\varepsilon} = (1 - V_{\varepsilon}^{a})(1 - V_{\varepsilon}^{b})f$,则 $g_{\varepsilon} \in A(\mathbf{T})$ 在 $[a, b]$ 的某邻域内为 $0$.又
\begin{align*}
& \|f - g_{\varepsilon}\|_{A(\mathbf{T})} \leq \|f - (1 - V_{\varepsilon}^{a})f\|_{A(\mathbf{T})} \\
& \qquad + \|(1 - V_{\varepsilon}^{a})f - (1 - V_{\varepsilon}^{a})(1 - V_{\varepsilon}^{b})f\|_{A(\mathbf{T})} \\
& \qquad \leq \|V_{\varepsilon}^{a}f\|_{A(\mathbf{T})} + \|1 - V_{\varepsilon}^{a}\|_{A(\mathbf{T})}\|V_{\varepsilon}^{b}f\|_{A(\mathbf{T})} \\
& \qquad \leq \|V_{\varepsilon}(f_{-a})\|_{A(\mathbf{T})} + 4\|V_{\varepsilon}(f_{-b})\|_{A(\mathbf{T})} \to 0 \quad (\varepsilon \to 0)
\end{align*}
故 $\|f - g_{\varepsilon}\|_{A(\mathbf{T})} \to 0$($\varepsilon \to 0$).
\end{proof}

\subsection{Malliavin 定理}\label{sec:3.3.2}

这里叙述交换 Banach 代数 $A(\mathbf{T})$ 谱综合不可能.今天,这一结果被称为 Malliavin 定理(cf.\ \cite{ch3:6}, \cite{ch3:7}, \cite{ch3:8}, \cite{ch3:11}, \cite{ch3:34}, \cite{ch3:39}).

\begin{theorem}[Malliavin]\label{thm:3.3.3}
$A(\mathbf{T})$ 谱综合不可能.
\end{theorem}

下面叙述的补题~\ref{lem:3.3.4},补题~\ref{lem:3.3.5} 证明上述定理.

\begin{lemma}\label{lem:3.3.4}
设实值函数 $f \in A(\mathbf{T})$ 具有下列性质:$t^{p}\omega(t) \in L^{1}(\mathbf{R})$($p = 0, 1, \ldots$),其中
\begin{equation*}
\omega(t) = \sup_{\|g\|_{A(\mathbf{T})} \leq 1}\left|\int e^{itf(x)}g(x)\frac{\dd x}{2\pi}\right|.
\end{equation*}
此时,存在 $a \in \mathbf{R}$ 使得由 $(f - a)^{p}$($p = 0, 1, \ldots$)生成的 $A(\mathbf{T})$ 的闭理想 $I_{p}$($p = 0, 1, \ldots$)全都不同.
\end{lemma}

\begin{proof}
令 $K = f(\mathbf{T})$,固定 $\tau \in C(\mathbf{T})$.考虑从 $\varphi \in C(K)$ 到 $\int \varphi(f(x))\tau(x)\frac{\dd x}{2\pi}$ 的复值函数,它是线性的,且
\begin{equation*}
\left|\int \varphi(f(x))\tau(x)\frac{\dd x}{2\pi}\right| \leq \|\varphi\|_{\infty}\|\tau\|_{\infty}\|\mu\|
\end{equation*}
故这是连续线性泛函.因此由 Riesz 表示定理,存在 $\lambda_{\tau} \in M(K)$($K$ 上有界正则 Borel 测度的全体)使得
\begin{equation*}
\int \varphi(f(x))\tau(x)\frac{\dd x}{2\pi} = \int_{K}\varphi(t)\dd\lambda_{\tau}(t).
\end{equation*}
这里,考察由 $\tau \in C(\mathbf{T})$ 到 $\lambda_{\tau} \in M(K)$ 的映射,线性是显然的.进一步,当 $\psi \in C(K)$ 时,$d\lambda_{\psi(f)\tau} = \psi(x)\dd\lambda_{\tau}$ 成立.事实上,
\begin{equation*}
\int_{K}\varphi(t)\psi(t)\dd\lambda_{\tau}(t) = \int \varphi(f(x))\psi(f(x))\tau(x)\frac{\dd x}{2\pi} = \int_{K}\varphi(t)\dd\lambda_{\psi(f)\tau}(t).
\end{equation*}
于是令 $\varphi(t) = e^{ist}$,$\tau = g(x)$,则
\begin{equation*}
\int e^{isf(x)}g(x)\frac{\dd x}{2\pi} = \int_{K}e^{ist}\dd\lambda_{g}(t).
\end{equation*}
此时,利用假设,
\begin{equation*}
\left|\int_{K}e^{ist}\dd\lambda_{g}\right| \leq \omega(s)\|g\|_{A(\mathbf{T})}, \quad s^{p}\omega(s) \in L^{1}(\mathbf{R}) \ (p = 0, 1, \ldots)
\end{equation*}
故可写为 $d\lambda_{g}(t) = m_{g}(t)\dd t$,$m_{g} \in C^{\infty}(\mathbf{R})$($\mathbf{R}$ 上无限次可微函数的全体).因此
\begin{equation*}
m_{g}(t) = \frac{1}{2\pi}\int_{-\infty}^{\infty}e^{its}\widehat{\lambda_{g}}(s)\dd s
\end{equation*}
得
\begin{align*}
\frac{\dd^{p}}{\dd t^{p}}m_{g}(t) & = \frac{1}{2\pi}\int_{-\infty}^{\infty}(is)^{p}e^{its}\widehat{\lambda_{g}}(s)\dd s \\
& = \frac{i^{p}}{2\pi}\int_{-\infty}^{\infty}s^{p}e^{its}\widehat{\lambda_{g}}(s)\dd s.
\end{align*}
从而
\begin{equation*}
\left|\frac{\dd^{p}}{\dd t^{p}}m_{g}(t)\right| \leq \frac{1}{2\pi}\int_{-\infty}^{\infty}|s|^{p}\omega(s)\dd s\,\|g\|_{A(\mathbf{T})}
\end{equation*}
得
\begin{equation*}
\left|\frac{\dd^{p}}{\dd t^{p}}m_{g}(t)\right| \leq C_{p}\|g\|_{A(\mathbf{T})}.
\end{equation*}
因此,对 $a \in \mathbf{R}$,
\begin{equation*}
\Phi_{p}:\ g \mapsto \left(\frac{\dd^{p}}{\dd t^{p}}m_{g}\right)(a)
\end{equation*}
是 $A(\mathbf{T})$ 上的连续线性泛函.于是令 $g \equiv 1$,则
\begin{equation*}
\int_{-\infty}^{\infty}m_{1}(t)\dd t = \widehat{\lambda_{1}}(0) = \int \frac{\dd x}{2\pi} = 1 \neq 0
\end{equation*}
故取 $a \in \mathbf{R}$ 使 $m_{1}(a) \neq 0$,则对 $p, q = 1, 2, \ldots$,由 $\lambda_{\tau}$ 的性质,
\begin{align*}
\Phi_{p}((f - a)^{q}) & = \left(\frac{\dd^{p}}{\dd t^{p}}m_{(f-a)^{q}}\right)(a) \\
& = \left(\frac{\dd^{p}}{\dd t^{p}}(t - a)^{q}m_{1}\right)(a).
\end{align*}
因此,若 $q > p$,则 $\Phi_{p}((f - a)^{q}) = 0$;若 $q = p$,则 $\Phi_{p}((f - a)^{p} \cdot 1) = p!\,m_{1}(a) \neq 0$,故对 $p = 1, 2, \ldots$,由 $(f - a)^{p}$ 生成的各个理想全都不同.
\end{proof}

其次,证明满足补题~\ref{lem:3.3.4} 假设的 $f$ 的存在性.

\begin{lemma}\label{lem:3.3.5}
存在实值函数 $f \in A(\mathbf{T})$ 与 $\delta > 0$ 满足
\begin{equation*}
\omega(t) = \sup_{\|g\|_{A(\mathbf{T})} \leq 1}\left|\int_{0}^{2\pi}e^{itf(x)}g(x)\dd x\right| = O\left(e^{-\delta|t|^{1/2}}\right).
\end{equation*}
由此可证明 $f$ 满足补题~\ref{lem:3.3.4}.
\end{lemma}

\begin{proof}
第 1 阶段.设 $0 \neq g_{\alpha}$,$0 \neq h_{\alpha} \in A(\mathbf{T})$($\alpha \in \mathbf{R}$)连续依赖于 $\alpha$.此时,对任意 $M > 0$,$\varepsilon > 0$,存在 $k \in \mathbf{N}$ 使得
\begin{equation*}
\|g_{\alpha}(x)h_{\alpha}(kx)\|_{A^{\infty}} \leq (1 + \varepsilon)\|g_{\alpha}\|_{A^{\infty}}\|h_{\alpha}\|_{A^{\infty}} \quad (|\alpha| \leq M).
\end{equation*}
其中 $\|f\|_{A^{\infty}} = \sup_{n \in \mathbf{Z}}\left|\hat{f}(n)\right|$.

第 1 阶段的证明.固定 $\alpha$,设 $k$ 为自然数.
\begin{equation*}
\widehat{g_{\alpha}(x)h_{\alpha}(kx)}(l) = \sum_{j}\widehat{h_{\alpha}}(j)\widehat{g_{\alpha}}(l - kj)
\end{equation*}
取 $l'$,$j'$ 满足 $l' = l - kj'$,$-k/2 \leq l' \leq (k/2) + 1$,则
\begin{equation*}
\left|\sum_{j}\widehat{g_{\alpha}}(l - kj)\widehat{h_{\alpha}}(j) - \widehat{g_{\alpha}}(l')\widehat{h_{\alpha}}(j')\right| \leq \sum_{|n| \geq \frac{k}{2} - 1}\left|\widehat{g_{\alpha}}(n)\right|\|h_{\alpha}\|_{A^{\infty}}.
\end{equation*}
因此
\begin{equation*}
\|g_{\alpha}(x)h_{\alpha}(kx)\|_{A^{\infty}} \leq \sum_{|n| \geq \frac{k}{2} - 1}\left|\widehat{g_{\alpha}}(n)\right|\|h_{\alpha}\|_{A^{\infty}} + \|g_{\alpha}\|_{A^{\infty}}\|h_{\alpha}\|_{A^{\infty}}
\end{equation*}
以及
\begin{equation*}
\left|\widehat{g_{\alpha}}(l')\widehat{h_{\alpha}}(j')\right| \leq \sum_{|n| \geq \frac{k}{2} - 1}\left|\widehat{g_{\alpha}}(n)\right|\|h_{\alpha}\|_{A^{\infty}} + \|g_{\alpha}(x)h_{\alpha}(kx)\|_{A^{\infty}}.
\end{equation*}
由后一式得
\begin{equation*}
\|g_{\alpha}\|_{A^{\infty}}\|h_{\alpha}\|_{A^{\infty}} \leq \sum_{|n| \geq \frac{k}{2} - 1}\left|\widehat{g_{\alpha}}(n)\right|\|h_{\alpha}\|_{A^{\infty}} + \|g_{\alpha}(x)h_{\alpha}(kx)\|_{A^{\infty}}.
\end{equation*}
因此,固定 $\alpha$ 时,对 $\varepsilon > 0$ 取充分大的 $k$,
\begin{equation*}
\left|\|g_{\alpha}\|_{A^{\infty}}\|h_{\alpha}\|_{A^{\infty}} - \|g_{\alpha}(x)h_{\alpha}(kx)\|_{A^{\infty}}\right| < \varepsilon.
\end{equation*}
这里,利用 $[-M, M]$ 是紧集,对 $\eta > 0$ 存在 $\delta > 0$ 与 $\alpha_{1}, \ldots, \alpha_{n} \in [-M, M]$,使得对任意 $|\alpha| \leq M$,存在满足 $|\alpha - \alpha_{j}| < \delta$ 的 $\alpha_{j}$,且 $\|g_{\alpha} - g_{\alpha_{j}}\|_{A(\mathbf{T})} < \eta$,$\|h_{\alpha} - h_{\alpha_{j}}\|_{A(\mathbf{T})} < \eta$.此时,取充分大的 $k$,
\begin{equation*}
\left|\|g_{\alpha_{j}}\|_{A^{\infty}}\|h_{\alpha_{j}}\|_{A^{\infty}} - \|g_{\alpha_{j}}(x)h_{\alpha_{j}}(kx)\|_{A^{\infty}}\right| < \eta \quad (j = 1, \ldots, n),
\end{equation*}
令
\begin{equation*}
C = \sup_{\alpha \in [-M, M]}\|g_{\alpha}\|_{A^{\infty}} < \infty, \quad D = \sup_{\alpha \in [-M, M]}\|h_{\alpha}\|_{A^{\infty}} < \infty,
\end{equation*}
则对 $\alpha \in [-M, M]$,存在 $\alpha_{j}$,
\begin{align*}
& \left|\|g_{\alpha}\|_{A^{\infty}}\|h_{\alpha}\|_{A^{\infty}} - \|g_{\alpha}(x)h_{\alpha}(kx)\|_{A^{\infty}}\right| \\
& \leq \left|\|g_{\alpha}\|_{A^{\infty}}\|h_{\alpha}\|_{A^{\infty}} - \|g_{\alpha_{j}}\|_{A^{\infty}}\|h_{\alpha_{j}}\|_{A^{\infty}}\right| \\
& \qquad + \left|\|g_{\alpha_{j}}\|_{A^{\infty}}\|h_{\alpha_{j}}\|_{A^{\infty}} - \|g_{\alpha_{j}}(x)h_{\alpha_{j}}(kx)\|_{A^{\infty}}\right| \\
& \qquad + \|g_{\alpha_{j}}(x)h_{\alpha_{j}}(kx) - g_{\alpha}(x)h_{\alpha_{j}}(kx)\|_{A^{\infty}} \\
& \qquad + \|g_{\alpha}(x)h_{\alpha_{j}}(kx) - g_{\alpha}(x)h_{\alpha}(kx)\|_{A^{\infty}} \\
& \leq \|g_{\alpha} - g_{\alpha_{j}}\|_{A(\mathbf{T})}\|h_{\alpha}\|_{A^{\infty}} + \|g_{\alpha_{j}}\|_{A^{\infty}}\|h_{\alpha} - h_{\alpha_{j}}\|_{A(\mathbf{T})} \\
& \qquad + \|g_{\alpha} - g_{\alpha_{j}}\|_{A(\mathbf{T})}\|h_{\alpha_{j}}\|_{A^{\infty}} + \|g_{\alpha}\|_{A^{\infty}}\|h_{\alpha} - h_{\alpha_{j}}\|_{A(\mathbf{T})} \\
& \leq 2(C + D)\eta.
\end{align*}
这里,把 $\eta > 0$ 取得足够小,使右边小于 $\varepsilon$ 即可.

第 2 阶段.对不含原点的闭区间 $[a, b]$,存在 $\rho$ 使得
\begin{equation*}
\left\|e^{i\alpha \cos kx}\right\|_{A^{\infty}} \leq \rho < 1 \quad (k \in \mathbf{Z},\ \alpha \in [a, b]).
\end{equation*}
由 $\|f(kx)\|_{A^{\infty}} = \|f(x)\|_{A^{\infty}}$,$\left\|e^{i\alpha \cos kx}\right\|_{A^{\infty}} = \|e^{i\alpha \cos x}\|_{A^{\infty}}$.又,令 $f_{\alpha}(x) = e^{i\alpha \cos x}$,利用 Parseval 等式,
\begin{equation*}
1 = \frac{1}{2\pi}\int_{0}^{2\pi}\left|e^{i\alpha \cos x}\right|^{2}\dd x = \sum_{n=-\infty}^{\infty}\left|\widehat{f_{\alpha}}(n)\right|^{2}
\end{equation*}
故对 $\alpha \in [a, b]$,由 $\alpha \neq 0$,$\left|\widehat{f_{\alpha}}(n)\right| < 1$($n \in \mathbf{Z}$).进一步,若 $\sup_{\alpha \in [a, b]}\sup_{n}\left|\widehat{f_{\alpha}}(n)\right| = 1$,则存在 $\{\alpha_{j}\}$,$\{n_{\alpha_{j}}\}$ 使 $\left|\widehat{f_{\alpha_{j}}}(n_{\alpha_{j}})\right| \to 1$($j \to \infty$).因 $[a, b]$ 是紧集,不妨设 $\alpha_{j} \to \alpha_{0}$($j \to \infty$).此时
\begin{equation*}
\left|\widehat{f_{\alpha_{j}}}(n) - \widehat{f_{\alpha_{0}}}(n)\right| \leq \int\left|e^{i\alpha_{j}\cos x} - e^{i\alpha_{0}\cos x}\right|\dd x \to 0 \quad (\alpha_{j} \to \alpha_{0})
\end{equation*}
且若 $|n_{\alpha_{j}}| \to \infty$,则 $\widehat{f_{\alpha_{0}}}(n_{\alpha_{j}}) \to 0$,故不妨设 $\{n_{\alpha_{j}}\}$ 有界.进一步,由 $\left|\widehat{f_{\alpha_{j}}}(n_{\alpha_{j}})\right| \to 1$($j \to \infty$),不妨设 $n_{\alpha_{j}} = n_{\alpha_{0}}$.即 $f_{\alpha_{0}}(x) = e^{in_{\alpha_{0}}x}$.这意味著 $e^{in_{\alpha_{0}}x} = e^{i\alpha\cos x}$($x \in \mathbf{T}$),产生矛盾.

第 3 阶段.令
\begin{equation*}
f(x) = \sum_{k=1}^{\infty}\frac{1}{k^{2}}\cos p_{k}x
\end{equation*}
时,适当选取 $\{p_{k}\} \subset \mathbf{N}$,则 $f(x)$ 即为所求.

第 3 阶段的证明.选取满足 $\prod_{k=1}^{\infty}(1 + \varepsilon_{k}) < 2$ 的 $\varepsilon_{k} > 0$,$p_{k}$ 如下选取.令 $p_{1} = 1$,设 $p_{n}$ 以前已选取.利用第 1 阶段选取 $p_{n+1}$:
\begin{align*}
& \left\|\prod_{k=1}^{n+1}e^{i(u/k^{2})\cos p_{k}x}\right\|_{A^{\infty}} \\
& = \left\|e^{i\sum_{k=1}^{n+1}(u/k^{2})\cos p_{k}x}\right\|_{A^{\infty}} \\
& \leq (1 + \varepsilon_{n})\left\|\prod_{k=1}^{n}e^{i(u/k^{2})\cos p_{k}x}\right\|_{A^{\infty}}\left\|e^{i(u/(n+1)^{2})\cos p_{n+1}x}\right\|_{A^{\infty}} \quad (|u| \leq n^{2}).
\end{align*}
这样选取 $\{p_{k}\} \subset \mathbf{N}$.此时,
\begin{equation*}
f(x) = \sum_{k=1}^{\infty}\frac{1}{k^{2}}\cos p_{k}x
\end{equation*}
即为所求.事实上,固定 $u \neq 0$.不妨设 $n \in \mathbf{N}$ 满足 $4^{n-1} \leq |u| < 4^{n}$.此时,由 $\left\|e^{i(u/k^{2})\cos p_{k}x}\right\|_{A^{\infty}} \leq 1$,
\begin{align*}
\|e^{iuf}\|_{A^{\infty}} & = \left\|e^{iu\sum_{k=1}^{\infty}(1/k^{2})\cos p_{k}x}\right\|_{A^{\infty}} \\
& \leq \left\|\prod_{k=1}^{2^{n+1}}e^{i(u/k^{2})\cos p_{k}x}\right\|_{A^{\infty}}\left\|\prod_{k=2^{n+1}+1}^{\infty}e^{i(u/k^{2})\cos p_{k}x}\right\|_{A^{\infty}} \\
& \leq 2\prod_{k=2^{n}}^{2^{n+1}}(1 + \varepsilon_{k})\left\|e^{i(u/k^{2})\cos p_{k}x}\right\|_{A^{\infty}}.
\end{align*}
这里,
\begin{equation*}
\frac{1}{16} = \frac{4^{n-1}}{4^{n+1}} \leq \frac{|u|}{k^{2}} < \frac{4^{n}}{4^{n}} = 1 \quad (2^{n} \leq k \leq 2^{n+1})
\end{equation*}
故 $1/16 \leq |u|/k^{2} < 1$.从而利用第 2 阶段($a = 1/16$,$b = 1$,$\alpha = |u|/k^{2}$),
\begin{equation*}
\|e^{iuf}\|_{A^{\infty}} \leq 2\left(\prod_{k=2^{n}}^{2^{n+1}}(1 + \varepsilon_{k})\right)\rho^{2^{n}} \quad (0 < \rho < 1).
\end{equation*}
因此,令 $\delta = -\log\rho$,
\begin{equation*}
\|e^{iuf}\|_{A^{\infty}} \leq 4e^{|u|^{1/2}\log\rho} = 4e^{-\delta|u|^{1/2}}.
\end{equation*}
于是,现在设 $g \in A(\mathbf{T})$,$\|g\|_{A(\mathbf{T})} \leq 1$,则
\begin{align*}
\left|\int_{0}^{2\pi}e^{itf(x)}g(x)\dd x\right| & = \left|\int_{0}^{2\pi}e^{itf(x)}\sum_{n \in \mathbf{Z}}\hat{g}(n)e^{inx}\dd x\right| \\
& \leq \sum|\hat{g}(n)|\left|\int_{0}^{2\pi}e^{itf(x)}e^{inx}\dd x\right| \\
& \leq 2\pi\|g\|_{A(\mathbf{T})}\|e^{itf}\|_{A^{\infty}}
\end{align*}
得 $\omega(t) = O\left(e^{-\delta|t|^{1/2}}\right)$.至此,第 3 阶段证明完毕.
\end{proof}

综上,命题得证.

\subsection{关于 Kronecker 集}\label{sec:3.3.3}

\begin{definition}\label{def:3.3.6}
$\mathbf{T}$ 的闭子集 $K$ 称为 Kronecker 集,是指对 $f \in S(K) = \{f \in C(K):\ |f(x)| = 1\ (x \in K)\}$,任意取 $\varepsilon > 0$ 时,存在 $n \in \mathbf{Z}$ 使得
\begin{equation*}
\max_{x \in K}|f(x) - e^{inx}| < \varepsilon.
\end{equation*}
\end{definition}

\begin{definition}\label{def:3.3.7}
$E \subset \mathbf{T}$ 在 $\mathbf{Q}$ 上线性无关,是指对 $\{x_{1}, \ldots, x_{k}\} \subset E \subset [-\pi, \pi) = \mathbf{T}$,$\{q_{0}, q_{1}, \ldots, q_{k}\} \subset \mathbf{Z}$,若 $q_{1}x_{1} + \cdots + q_{k}x_{k} = \pi q_{0}$,则 $q_{0} = q_{1} = \cdots = q_{k} = 0$.
\end{definition}

\textbf{Kronecker 定理}(\cite{ch3:8})．若 $\{x_{1}, x_{2}, \ldots, x_{n}\} \subset \mathbf{T}$ 在 $\mathbf{Q}$ 上线性无关,则对任意 $\{\beta_{1}, \beta_{2}, \ldots, \beta_{n}\} \subset \mathbf{R}$ 与任意 $\varepsilon > 0$,存在 $k \in \mathbf{Z}$ 使得
\begin{equation*}
\left|e^{ikx_{j}} - e^{i\beta_{j}}\right| < \varepsilon \quad (j = 1, \ldots, n).
\end{equation*}

下面,利用 Kronecker 定理容易证明.

\begin{theorem}\label{thm:3.3.8}
$\mathbf{T}$ 的线性无关的有限子集 $E$ 是 Kronecker 集.
\end{theorem}

\begin{theorem}\label{thm:3.3.9}
Kronecker 集 $E$ 是线性无关的集合.
\end{theorem}

\begin{proof}
对 $\{x_{1}, \ldots, x_{k}\} \subset E$,$q_{j} \in \mathbf{Q}$($j = 0, 1, \ldots, k$),设
\begin{equation*}
\sum_{j=1}^{k}q_{j}x_{j} = 2\pi q_{0}.
\end{equation*}
此时,对 $n \in \mathbf{Z}$,
\begin{equation*}
\prod_{j=1}^{k}(e^{inx_{j}})^{q_{j}} = e^{in\sum_{j=1}^{k}q_{j}x_{j}} = e^{i2\pi q_{0}n}
\end{equation*}
因此 $\prod_{j=1}^{k}(e^{inx_{j}})^{q_{j}} = 1$.这里,由 Kronecker 集的定义,取 $\alpha_{j} \in \mathbf{R}$($j = 1, \ldots, k$),则对 $\varepsilon > 0$ 存在 $m \in \mathbf{Z}$ 使得 $\left|e^{imx_{j}} - e^{i\alpha_{j}}\right| < \varepsilon$($j = 1, \ldots, k$).因此 $e^{i\sum_{j=1}^{k}\alpha_{j}q_{j}} = 1$($\alpha_{j} \in \mathbf{R}$).因此 $q_{j} = 0$($j = 1, \ldots, k$).从而 $q_{0} = 0$.综上,$E$ 线性无关.
\end{proof}

下面,由 Rudin \cite{ch3:11} 证明非平凡完全不连通完全集(称为 Cantor 集)构成的 Kronecker 集的存在性.

\begin{lemma}\label{lem:3.3.10}
设 $U_{1}, U_{2}, \ldots, U_{k}$ 为 $(0, \pi)$ 的非空互不相交开区间列,则存在 $x_{j} \in U_{j}$($j = 1, \ldots, k$)使 $\{x_{1}, \ldots, x_{k}\}$ 是 Kronecker 集.
\end{lemma}

\begin{proof}
由于 $U_{j}$ 是不可数集,可取 $x_{j} \in U_{j}$($j = 1, \ldots, k$)使得 $\{x_{1}, \ldots, x_{k}\}$ 线性无关.此时,由定理~\ref{thm:3.3.8} 成立.
\end{proof}

\begin{theorem}\label{thm:3.3.11}
$\mathbf{T}$ 中存在是 Cantor 集的 Kronecker 集.
\end{theorem}

\begin{proof}
令 $P_{1}^{0} = \left[\frac{1}{3}, \frac{1}{2}\right] \subset (-\pi, \pi)$.设对 $r \geq 1$,已作出 $s = 2^{r-1}$ 个互不相交的闭区间 $P_{1}^{r-1}, P_{2}^{r-1}, \ldots, P_{s}^{r-1}$.此时,在 $P_{j}^{r-1}$ 中取互不相交的开区间 $W_{2j-1}, W_{2j}$($1 \leq j \leq s$).这样,由补题~\ref{lem:3.3.10},存在 Kronecker 集 $\{x_{1}^{r}, \ldots, x_{2s}^{r}\}$($x_{j}^{r} \in W_{j}$,$1 \leq j \leq 2s$).这里,存在 $A_{r} \subset \mathbf{Z}$(有限集)使得对任意 $\{\alpha_{1}, \ldots, \alpha_{2^{r}}\} \subset \mathbf{R}$,取适当的 $n \in A_{r}$,
\begin{equation*}
\left|e^{i\alpha_{j}} - e^{inx_{j}^{r}}\right| < \frac{1}{r} \quad (j = 1, \ldots, 2^{r}).
\end{equation*}
此时,存在闭区间 $P_{j}^{r} \subset W_{j}$ 使得
\begin{equation*}
\left|e^{inx} - e^{inx_{j}^{r}}\right| < \frac{1}{r} \quad (x \in P_{j}^{r},\ n \in A_{r}).
\end{equation*}
进一步,设 $P_{j}^{r}$ 的长度 $\left|P_{j}^{r}\right| < 1/r$.如此归纳继续,令
\begin{equation*}
P = \bigcap_{r=1}^{\infty}\bigcup_{j=1}^{2^{r}}P_{j}^{r}.
\end{equation*}
由构造方法,$P$ 是 Cantor 集.设 $f \in S(P)$,$\varepsilon > 0$.由 $f \in C(P)$,$f$ 在 $P$ 上一致连续.因此存在足够大的 $r_{0}$ 使得 $f(P \cap P_{j}^{r_{0}})$($1 \leq j \leq 2^{r_{0}}$)不成为 $\{z \in \mathbf{C}:\ |z| = 1\}$ 的全部.此时,由 $P \cap P_{j}^{r_{0}} \subset (-\pi, \pi)$,存在 $\widetilde{f} \in S\left(\bigcup_{j}P_{j}^{r_{0}}\right)$ 使 $\widetilde{f}\Big|_{P \cap \left(\bigcup_{j}P_{j}^{r_{0}}\right)} = f$.进一步,由 $P_{j}^{r}$ 的构造方法,对 $\varepsilon > 0$,$r > r_{0}$,$r > 3/\varepsilon$,可设 $\left|\widetilde{f}(x) - \widetilde{f}(x_{j}^{r})\right| < \varepsilon/3$($x \in P_{j}^{r}$,$1 \leq j \leq 2^{r}$).又由构造方法,存在 $n \in A_{r}$ 使 $\left|\widetilde{f}(x_{j}^{r}) - e^{inx_{j}^{r}}\right| < 1/r$($1 \leq j \leq 2^{r}$).因此,对 $x \in \bigcup_{j}P_{j}^{r}$,
\begin{align*}
\left|\widetilde{f}(x) - e^{inx}\right| & \leq \left|\widetilde{f}(x) - \widetilde{f}(x_{j}^{r})\right| + \left|\widetilde{f}(x_{j}^{r}) - e^{inx_{j}^{r}}\right| + \left|e^{inx_{j}^{r}} - e^{inx}\right| \\
& < \frac{\varepsilon}{3} + \frac{1}{r} + \frac{1}{r} < \varepsilon.
\end{align*}
因此,$\left|f(x) - e^{inx}\right| < \varepsilon$($x \in P$).
\end{proof}

\begin{note}\label{note:3.3.12}
(cf.\ \cite{ch3:11})上述中令 $Q = P \cup (-P)$,取满足 $Q$ 上任意一点测度为 $0$ 的 $\mu \in M(\mathbf{T})$ 时,测度 $\delta_{0}, \mu, \mu^{2}, \mu^{3}, \ldots, \mu^{n} = \mu^{n-1} * \mu, \ldots$ 聚集在互不相交的集合上.这一点可由 $P$ 是线性无关集与测度乘积的定义证明.事实上,取 $n, m \geq 1$,令 $Q_{m} = \{x_{1} + \cdots + x_{m}:\ x_{j} \in Q\}$,则 $\mu^{n}(Q_{m}) = 0$($m < n$).例如,$n = 3$,$m = 1$,$\mu \geq 0$ 时,
\begin{equation}\tag{3.3.1}\begin{aligned}
\mu^{3}(Q_{m}) & = \displaystyle\int \mu^{2}(Q - x)\dd\mu(x) \\
& = \displaystyle\int\!\!\int \chi_{Q-x}(w)\dd\mu^{2}(w)\dd\mu(x) \\
& = \mu \times \mu \times \mu\bigl(\{(x, y, z) \in Q \times Q \times Q:\ x + y + z \in Q\}\bigr) \\
& \leq \sum_{i \neq j,\ 1 \leq i, j \leq 3}\mu \times \mu \times \mu\bigl(\{(x_{1}, x_{2}, x_{3}) \in Q \times Q \times Q:\ x_{i} \pm x_{j} = 0\}\bigr).
\end{aligned}\label{eq:3.3.1}\end{equation}
其中,$\chi_{Q-x}(w)$ 是 $Q - x$ 的特征函数.这里,用 Fubini 定理与 $\mu(\{x\}) = 0$($x \in \mathbf{T}$),\eqref{eq:3.3.1} 的右边为 $0$.用这样的想法可得 $\mu^{n}(Q_{m}) = 0$($m < n$).
\end{note}

由定义与 Riemann-Lebesgue 引理可知,Kronecker 集的 Lebesgue 测度为 $0$,且是完全不连通集.Varopoulos \cite{ch3:50} 利用这一性质对谱综合证明了下列结论.其证明是考察 $A(\mathbf{T})$ 上的连续线性泛函,充分利用 Kronecker 集的性质而进行的函数分析方法的证明,饶有趣味,但此处从略.证明请参照论文 \cite{ch3:50}.

\begin{theorem}\label{thm:3.3.13}
$\mathbf{T}$ 的 Kronecker 集是 $A(\mathbf{T})$ 的谱综合集.
\end{theorem}

即使在非离散局部紧 Abel 群上,也可用特征定义 Kronecker 集,但众所周知,一般未必完全不连通.Saeki \cite{ch3:44} 发展了 Varopoulos \cite{ch3:50} 的方法,得到了推广到非离散局部紧 Abel 群的结果.
